\section{Source parallelity in all handle dimensions and a bounded inherited core}
\label{sec:parallelity-core}

The full raw cut of Section~\ref{sec:cut-cells} retains the source geometry
needed to distinguish product regions from the exceptional pieces of a cut.
The new \code{normal\_cut\_parallelity} uses that structure to describe the
parallelity bundle for the supplied normal surface in the triangulation's
dual handle structure. It includes product handles of indices zero, one and
two, their complete base surfaces, their fibre twisting, and their interfaces
with an inherited complementary core. The number of complementary handles
is at most $30T$, independently of normal multiplicity, where $T$ is the
number of source tetrahedra.

The input is still a validated finite connected compact orientable
triangulation with one torus boundary and an admissible integral standard
normal vector. Empty, disconnected and one-sided surfaces are supported.
The output is an intermediate decomposition of the supplied cut. An
admissible patterned hierarchy and its global running-time bound still
require further work.

\subsection{Selecting the product handles}

Use the dual handle decomposition relative to the original boundary:
tetrahedron chambers give zero-handles, regions of internal faces give
one-handles, gaps on interior edges give two-handles, and interior original
vertices give three-handles. Boundary faces and boundary edges do not supply
one- or two-handles in this decomposition.

In a tetrahedron, a stack of $m$ identical normal discs has $m-1$ product
zero-handles between successive discs. Distinct disc types are separated by
exceptional chambers. On a canonical internal face, the arm at corner $v$
contains $A=t_v+q_{fv}$ parallel normal arcs. Its ranks $1,\ldots,A-1$ give
product one-handles; rank zero and the face centre are exceptional. Finally,
on an interior source edge of weight $w$, ranks $1,\ldots,w-1$ give product
two-handles. End gaps are exceptional. No three-handle is selected.

These selections are local and respect the source face permutations. At an
intersection of selected handles the product fibres agree, with an explicit
bit recording whether the local fibre parameter reverses. Their union is
therefore the source parallelity $I$-bundle. The rank-map slope is retained
separately from the fibre bit: an interval consisting of a single point
cannot encode a fibre reversal by its point map alone.

\subsection{Complete base polygons, including contacts with the core}

Each selected zero-handle contributes a truncated normal polygon: a hexagon
for a triangle and an octagon for a quadrilateral. Its sides alternate
between contacts across faces and contacts around source edges. A selected
one-handle contributes a rectangle, with two tetrahedron ends and two edge
sides. A selected two-handle contributes its edge-link polygon, alternating
tetrahedron sectors and internal-face sectors. The source edge link records
every occurrence, including repeated tetrahedra and self-face pairings.

All zero--one, zero--two and one--two contacts are glued along their actual
base arcs. Each contact has an affine rank map, a permutation of its arc
endpoints, and a fibre-reversal bit. Unglued arcs either lie on the original
boundary or carry a source-labelled port to a complementary zero- or
one-handle. A missing selected lower-index handle leaves a free boundary arc;
it is not silently filled by a disc.

This last point is essential. The discussion of boundary handles in the
retained July manuscript~\cite{lackenby2026}, in the proof of its lemma on
cutting along a vertical surface, explains why the naive list of selected
zero-, one- and two-handles need not itself be a handle decomposition of the
base. Higher-index handles can meet the complementary core. The polygon
construction here retains those boundary contacts explicitly. It implements
the source-cut parallelity geometry relevant to the cutting theorem, while
the later interior parallelity bundle of an admissible patterned structure
has different input hypotheses. The published algorithmic precedent is
Lackenby's full parallelity construction~\cite{lackenby-parallelity}.

A concrete case in the audit is the two-tetrahedron source labelled
\code{fibonacci\_lst\_02}, with coordinate rows
$(0,0,1,1,1,0,0)$ and $(1,1,1,1,0,0,0)$. Its parallelity bundle has one
component whose base is a disc, supplied by a single selected two-handle.
It has no selected zero- or one-handles and has eight core-facing arcs.
A prism-only extraction would miss this component entirely. Its attachment
subset is the whole base boundary, so $\chi(A)=0$, although all eight arcs
and their cyclic endpoint identifications must still be retained.

\subsection{The complementary handles and attachment charts}

\begin{proposition}
The unselected handles form an inherited complementary handle structure
with at most $6T$ zero-handles, $8T$ one-handles, $12T$ two-handles and $4T$
three-handles. Its attaching data and its interfaces with the bundle have
$O(T)$ source records with binary ranks.
\end{proposition}
\begin{proof}
Each tetrahedron has at most four triangle tips and two endpoints of an
occupied quadrilateral chain. With no quadrilateral stack the single centre
replaces those two endpoints. Hence there are at most six exceptional
chambers. Each internal face has its centre and at most three nonempty arm
tips, and there are at most $2T$ paired faces. Each interior edge has at most
two end gaps; the total number of source edges is at most $6T$. There are at
most $4T$ source vertices. Adding gives the stated $30T$ bound.

A complementary face region has only complementary chamber ends. An end
gap around an interior edge similarly meets only complementary chamber and
face regions. The source vertex links provide the three-handle charts.
Thus the lower attachments of complementary handles remain in the core.
The total length of the edge-link sector lists is bounded by the number of
local tetrahedron-edge flags; the vertex-link triangles are likewise bounded
by local vertex flags. Retaining each flag with its source chart therefore
uses linear combinatorial storage.
\end{proof}

The packet contains both ends of every core one-handle, the alternating
zero/one attaching sectors of every core two-handle, and the source
vertex-link triangles of every core three-handle. It also includes a copied
source triangulation and the unscaled normal coordinates. This preserves
topological attaching information beyond a cellular boundary matrix.

Every core-facing base arc is at a type transition and has singleton rank.
A rectangle end has at most one such transition per face corner and side.
An edge-cap tetrahedron sector has at most two, and a face sector has at
most one. These yield a linear port bound; the implementation conservatively
checks $66T$. Adjacent free cap arcs can face different core handles, even
when their integer ranks are consecutive. The constructor splits them at
the source triangle/quad transitions before assigning owners.

\subsection{Classifying the bases without expanding polygons}

Four interval-orbit queries suffice. First identify polygon corners across
glued arc endpoints. Next add the free-arc endpoint pairings: among these
new orbits, those containing free arcs are the boundary loops. The weighted
selectors of Section~\ref{sec:weighted-transversal} distinguish them from
isolated interior-vertex orbits and place exactly one representative mark
on each actual loop.

For each connected polygon component accumulate the seven coordinates
\[
 (\chi,b,h_0,h_1,h_2,e_{\mathrm{old}},m).
\]
Here $b$ is the number of base boundary components, $h_i$ counts selected
source handles of index $i$, $e_{\mathrm{old}}$ counts free arcs on the
original boundary, and $m$ is the bit mask of core ports. A polygon of
degree $d$ initially contributes $1-d$ to $\chi$. Each glued arc contributes
one back, because it was counted twice, and one mark at each quotient vertex
completes $V-E+F$. Boundary-loop marks contribute to $b$. Unique port $j$
contributes $2^j$ to $m$, so its encoded bit length is $O(T)$.

The fourth query is the fibre double cover of the polygon system. Let $C(w)$
be the base-component histogram and $D(w)$ its doubled histogram. If $O(w)$
and $N(w)$ count product and twisted components, respectively, then
\[
 C(w)=O(w)+N(w),\qquad D(w)=2O(w)+N(w/2),
\]
where the second term is zero when any coordinate of $w$ is odd. Process
signatures in increasing positive $h_0+h_1+h_2$. This recovers both types
without confusing an annulus signature with twice a M\"obius-band signature.
Since the ambient manifold is orientable, the total interval bundle is a
product exactly when its base is orientable. The genus or crosscap count
then follows from $\chi=2-2g-b$ or $\chi=2-k-b$.

\subsection{Checking the gluing to the core}

Let $A\subseteq\partial B$ be the union of core-labelled arcs in all base
surfaces, and let $K$ be the complementary core. The bundle meets the core
in the interval bundle over $A$. The vertex selector also records whether a
vertex is incident to a core port, giving
\[
 \chi(A)=\#\{\text{vertices incident to core ports}\}-\#\{\text{core ports}\}.
\]
The checked identities are
\[
 \chi(M\mathbin{\backslash}\operatorname{int}\nu(F))
   =\chi(K)+\chi(B)-\chi(A),\qquad
 \chi(B)-\chi(A)=h_0-h_1+h_2.
\]
The second identity explains why counting selected handles alone can give
the wrong Euler characteristic of the base. Every port belongs to exactly
one base component. A profile containing a port must have multiplicity one;
only components with no core contacts can be grouped into repeated profiles.
The source polygon packet retains their attachment locations even when
profiles alone omit the cyclic order of free arcs.

Independent verification reconstructs the source geometry from the raw
cut-cell reference model. It derives contact maps by inverting
edge-to-face-to-chamber incidences, rather than calling the producer's direct
rank formulas. It independently assembles all four query systems, checks the
two tagged selector proofs and the two weighted proofs, and recomputes the
surface classification and gluing identities. A shared discovery-cycle cap
covers all four queries; exhausted work has no geometric result. The
verification operation cap counts the supplied traces, and cancellation
propagates through both paths.

\subsection{Encoded size, experiments and remaining obligations}

There are at most $17T$ polygon families, $44T$ gluing families and $68T$
polygon side slots. If $D$ is the total number of normal discs, the polygon
and corner universes have sizes at most $9D$ and $32D$, respectively. These
universes are represented by binary interval endpoints. Their storage does
not scale with their numerical sizes. The source geometry and core have
$O(T)$ records; the attached weights and traces are additional encoded data.
As with the existing weighted engine, the implementation's replay accounting
is polynomial in the realized trace and weighted partitions. The present
work does not turn that conditional accounting into a complete hierarchy
complexity theorem.

The finite reference expands the polygons, checks each vertex link is a
single interval or circle, checks fibre and base orientation independently,
and counts connected components, boundary loops, ports and Euler
characteristics. The accompanying audit records its source hashes and
case-by-case proof hashes. The regression tests also cover one-sided
multiples, interior three-handles, relabellings, forged fibre and endpoint
maps, forged core owners, shared caps, cancellation, and exact JSON replay
of $2^{16384}$ parallel meridians with a fixed four-zero/three-one core.

\input{parallelity_results.tex}

There is also a normality obligation at the interface to the hierarchy.
The accepted standard coordinates satisfy triangulation matching and
quadrilateral constraints. The patterned handle-normality definition in
\cite{lackenby2026} imposes further local restrictions involving the
characteristic surface and boundary pattern. The manuscript's nonnegative
index lemma assumes that stronger normality and an admissible input handle
structure. A standard coordinate vector and a bounded source core alone
therefore do not authorize applying that lemma. The patterned construction
must supply and check these hypotheses explicitly.

The empty cut on the one-tetrahedron layered solid torus is already a
concrete obstruction to assuming admissibility. Its independently replayed
core has handle counts $(1,1,0,0)$. The characteristic surface has two
zero-handles, the ends of the single manifold one-handle, and no one-handles,
since there are no manifold two-handles. With empty pattern, each
characteristic zero-handle has index $-4+2\cdot0+0=-4$. Thus this inherited
core fails the first admissibility condition. The saved source certificate
in \code{parallelity-admissibility-obstruction.json} makes the distinction
checkable even before any large normal vector is introduced.

This operator supplies the all-dimensional source parallelity extraction
missing from the earlier prism-only construction. The next requirements are
to carry boundary patterns through the core interfaces, construct and verify
an admissible uniform-type structure, and implement the hierarchy operations
with their reduction and search bounds. The complete recognizer's general
exponential fallback remains. Quasi-polynomial unknot recognition has not
yet been established by this implementation.
